\section{Madelung Hydrodynamics}

The Madelung representation writes the wavefunction as
\begin{equation}
\psi=\sqrt{\rho}\,e^{iS/\hbar}.
\end{equation}

Substitution into the Schr\"odinger equation yields:
\begin{equation}
\partial_t \rho+\nabla\cdot\left(\rho\frac{\nabla S}{m}\right)=0,
\label{continuity}
\end{equation}
and
\begin{equation}
\partial_t S+\frac{(\nabla S)^2}{2m}+V
    -\frac{\hbar^2}{2m}\frac{\Delta\sqrt{\rho}}{\sqrt{\rho}}=0.
\label{madelungHJ}
\end{equation}

At first sight this resembles ordinary hydrodynamics. However, the velocity field
\[v=\frac{\nabla S}{m}\]
is highly constrained.

Away from singularities,
\begin{equation}
\nabla\times v=0.
\end{equation}
Thus Madelung fluids are not generic fluids but irrotational potential flows.

\section{Multivalued Phases and Interference}

After caustic formation, semiclassical quantum mechanics produces
\begin{equation}
\psi \sim \sum_j A_j e^{iS_j/\hbar}.
\end{equation}
Thus quantum mechanics naturally accommodates multiple classical branches simultaneously.

The hydrodynamic velocity becomes
\begin{equation}
v=\frac{\hbar}{m}\operatorname{Im}\frac{\nabla\psi}{\psi},
\end{equation}
which generally is not equal to the gradient of a single classical action branch.
Near nodes where $\psi=0$,
the velocity field becomes singular and vortices appear.

\section{Global Constraints and Quantization}

Single-valuedness of the wavefunction implies 
\begin{equation} 
    \oint \nabla S\cdot dl =2\pi n\hbar. 
\end{equation} 
Thus the Madelung variables are globally
constrained. Without imposing these quantization conditions, the Madelung
system admits more solutions than Schr\"odinger theory. This issue is closely
related to the Wallstrom objection in stochastic mechanics.

\section{Consequences for Nelson Mechanics}

Nelson stochastic mechanics introduces a current velocity
\[v=\frac{\nabla S}{m}\]
together with stochastic diffusion.

However, unless additional global constraints are imposed, the resulting theory
generally possesses too many states compared with quantum mechanics.

The underlying issue is again the hidden geometric structure inherited from:
\begin{itemize}
\item the original two-point action;
\item the multivalued semiclassical structure;
\item global phase coherence.
\end{itemize}

\section{Conclusion}

The Hamilton--Jacobi equation does not truly replace a two-point dynamical
object by a one-point field. The apparent scalar field $S(q,t)$ is only a local
representation of an evolving Lagrangian manifold in phase space.

The gradient representation $p=\nabla S$ fails generically after caustic
formation, when the phase-space manifold folds and multiple classical
trajectories reach the same configuration point.
Madelung hydrodynamics inherits these limitations. The velocity field is not
globally well-defined, and the true quantum object remains the wavefunction
itself rather than a classical fluid velocity field.

Thus the hidden two-point structure of classical action survives in:
\begin{itemize}
\item the geometry of Lagrangian manifolds,
\item multivalued semiclassical phases,
\item quantization conditions,
\item and the global topology of wave mechanics.
\end{itemize}
